\section{Problema}
Una estación de radio vende un bloque publicitario de duración fija, dividido en pautas de duración entera. Hay una tarifa fija a la que se vende cada duración posible.

Se quiere partir el bloque completo en pautas, cada una de al menos una unidad de duración, para maximizar el ingreso total. Diseñe un algoritmo de programación dinámica para ello.

\begin{itemize}
  \item Entrada: Un arreglo \inline{tarifa} donde \inline{tarifa[i]} es el precio de una pauta de duración $i$.
  \item Salida: El ingreso máximo obtenible al partir un bloque de duración $n$ en pautas.
\end{itemize}

E.g., para \inline{tarifa = [0, 3, 5, 8, 9]}, el mejor reparto de un bloque de duración $4$ es en cuatro pautas de duración $1$, con un ingreso de $12$.

\section{Entregables}

\begin{enumerate}
  \item (2 puntos) Defina una relación de recurrencia que resuelve el problema. Explique en palabras qué significa la función y qué significa cada variable. Sea sucinto.
  \item (1 punto) ¿Cuánto espacio requiere una solución de programación dinámica para este problema? ¿Se puede optimizar? Justifique su respuesta con algún modelo matemático del problema o de la solución.
  \item (2 puntos) Implemente una solución top-down o bottom-up (a su elección) en Python, Java, C, C++, JavaScript o TypeScript.
\end{enumerate}

\section{Rúbrica}

\begin{itemize}
  \item Si declara no saber responder un numeral, se acredita automáticamente el 20\% de su valor. Puede escribir: ``No sé''.
  \item Si deja un numeral en blanco, el siguiente asume su valor. E.g., si el numeral 1 queda en blanco, el numeral 2 vale 3 puntos.
\end{itemize}
